\documentclass[11pt]{article}
\usepackage{amsmath,amssymb,amsthm}
\usepackage{algorithm}
\usepackage{algorithmic}
\usepackage{enumitem}
\usepackage{hyperref}
\usepackage{bm}
\usepackage{geometry}
\geometry{margin=1in}
\newtheorem{theorem}{Theorem}[section]
\newtheorem{lemma}{Lemma}[section]
\newtheorem{assumption}{Assumption}[section]
\newtheorem{corollary}{Corollary}[section]

\begin{document}

\title{Randomized Block Coordinate Descent (RBCD)}
\author{}
\date{}
\maketitle

%%%%%%%%%%%%%%%%%%%%%%%%%%%%%%%%%%%%%%%%%%%%%%%%%%%%%%%%%%%%%%%%%%%%%%%%%%%%%%%
\section*{Algorithm Details Expansion}
%%%%%%%%%%%%%%%%%%%%%%%%%%%%%%%%%%%%%%%%%%%%%%%%%%%%%%%%%%%%%%%%%%%%%%%%%%%%%%%

\section{Mathematical Formulation}
Let $f:\mathbb{R}^{d}\rightarrow\mathbb{R}$ be a convex differentiable function.
The vector $w\in\mathbb{R}^{d}$ is partitioned into $B$ (non‑overlapping) blocks
\[
w=\bigl(w^{(1)},w^{(2)},\dots ,w^{(B)}\bigr),\qquad 
w^{(i)}\in\mathbb{R}^{d_i},\;\;\sum_{i=1}^{B}d_i=d .
\]

At iteration $t\ge 0$ we draw a block index $i_t\in\{1,\dots ,B\}$ according to a
probability distribution $p=(p_1,\dots ,p_B)$ (in the sequel we take the
uniform distribution $p_i=1/B$).  
Define the *partial gradient* with respect to block $i$ as
\[
\nabla_i f(w) \;:=\; \bigl[\,\partial_{w^{(i)}}f(w)\,\bigr]\in\mathbb{R}^{d_i}.
\]

Given a stepsize $\eta_t>0$ (possibly block‑dependent), the update reads
\[
\boxed{
\begin{aligned}
w_{t+1}^{(i)} &=
\begin{cases}
w_{t}^{(i)}-\eta_t^{(i)}\,\nabla_i f(w_t), & \text{if } i=i_t,\\[4pt]
w_{t}^{(i)}, & \text{otherwise},
\end{cases}\\[4pt]
w_{t+1} &:=\bigl(w_{t+1}^{(1)},\dots ,w_{t+1}^{(B)}\bigr).
\end{aligned}}
\tag{1}
\]

Typical choices for the stepsize are the *block‑Lipschitz* choice
\[
\eta_t^{(i)}=\frac{1}{L_i}, \qquad
L_i\ \text{such that}\ 
\|\nabla_i f(u)-\nabla_i f(v)\| \le L_i\|u^{(i)}-v^{(i)}\|
\ \forall u,v\in\mathbb{R}^d .
\tag{2}
\]

%%%%%%%%%%%%%%%%%%%%%%%%%%%%%%%%%%%%%%%%%%%%%%%%%%%%%%%%%%%%%%%%%%%%%%%%%%%%%%%
\section{Pseudocode}
%%%%%%%%%%%%%%%%%%%%%%%%%%%%%%%%%%%%%%%%%%%%%%%%%%%%%%%%%%%%%%%%%%%%%%%%%%%%%%%

\begin{algorithm}[H]
\caption{Randomized Block Coordinate Descent (RBCD)}
\label{alg:RBCD}
\begin{algorithmic}[1]
\REQUIRE Initial point $w_0\in\mathbb{R}^d$, stepsizes $\{\eta^{(i)}\}_{i=1}^B$, 
number of iterations $T$.
\FOR{$t=0,\dots,T-1$}
    \STATE Sample a block $i_t\in\{1,\dots ,B\}$ with probability $p_{i_t}=1/B$.
    \STATE Compute the partial gradient $\displaystyle g_t:=\nabla_{i_t}f(w_t)$.
    \STATE Update the selected block:
    \[
    w_{t+1}^{(i_t)} \gets w_{t}^{(i_t)}-\eta^{(i_t)} g_t .
    \]
    \STATE Keep all other blocks unchanged:
    \[
    w_{t+1}^{(i)} \gets w_{t}^{(i)},\qquad\forall i\neq i_t .
    \]
\ENDFOR
\RETURN $w_T$ (or optionally the averaged iterate $\bar w_T:=\frac1T\sum_{t=0}^{T-1}w_t$).
\end{algorithmic}
\end{algorithm}

%%%%%%%%%%%%%%%%%%%%%%%%%%%%%%%%%%%%%%%%%%%%%%%%%%%%%%%%%%%%%%%%%%%%%%%%%%%%%%%
\section{Dry Run Example}
%%%%%%%%%%%%%%%%%%%%%%%%%%%%%%%%%%%%%%%%%%%%%%%%%%%%%%%%%%%%%%%%%%%%%%%%%%%%%%%

We illustrate the algorithm on a one‑dimensional problem (hence a single
block).  
\[
f(w)=(w-3)^2,\qquad w_0=0,\qquad \eta=0.1 .
\]

Because $B=1$, the block index $i_t$ is always $1$.  
The full gradient equals the partial gradient:
\[
\nabla f(w)=2(w-3).
\]

\begin{center}
\begin{tabular}{c|c|c|c}
Iteration $t$ & $w_t$ & $\nabla f(w_t)$ & $w_{t+1}=w_t-\eta\nabla f(w_t)$\\\hline
0 & $0$ & $2(0-3)=-6$ & $0-0.1(-6)=0.6$\\
1 & $0.6$ & $2(0.6-3)=-4.8$ & $0.6-0.1(-4.8)=1.08$\\
2 & $1.08$ & $2(1.08-3)=-3.84$ & $1.08-0.1(-3.84)=1.464$\\\hline
\end{tabular}
\end{center}

Corresponding objective values:

\[
f(0)=9,\; f(0.6)=5.76,\; f(1.08)=3.6864,\; f(1.464)=2.348.
\]

The iterates move steadily towards the minimiser $w^\star=3$.

%%%%%%%%%%%%%%%%%%%%%%%%%%%%%%%%%%%%%%%%%%%%%%%%%%%%%%%%%%%%%%%%%%%%%%%%%%%%%%%
\section*{Convergence Theory Development}
%%%%%%%%%%%%%%%%%%%%%%%%%%%%%%%%%%%%%%%%%%%%%%%%%%%%%%%%%%%%%%%%%%%%%%%%%%%%%%%

\section{Assumptions}
\begin{assumption}[Convexity]\label{as:convex}
$f:\mathbb{R}^d\to\mathbb{R}$ is convex, i.e.
\[
f(\theta u+(1-\theta)v)\le \theta f(u)+(1-\theta)f(v),\quad\forall u,v\in\mathbb{R}^d,\;\theta\in[0,1].
\]
\end{assumption}

\begin{assumption}[Block‑Coordinate Lipschitz Smoothness]\label{as:lipschitz}
For each block $i\in\{1,\dots ,B\}$ there exists $L_i>0$ such that for all
$u,v\in\mathbb{R}^d$,
\[
\bigl\|\nabla_i f(u)-\nabla_i f(v)\bigr\|
\le L_i\bigl\|u^{(i)}-v^{(i)}\bigr\|.
\tag{3}
\]
\end{assumption}

\begin{assumption}[Bounded Level Set]\label{as:bounded}
Let $w^\star$ be any minimiser of $f$. The initial point $w_0$ satisfies
\[
\|w_0-w^\star\|_{(L)}^2:=\sum_{i=1}^B L_i\|w_0^{(i)}-w^{\star (i)}\|^2<\infty .
\tag{4}
\]
\end{assumption}

\begin{assumption}[Sampling]\label{as:sampling}
Blocks are sampled i.i.d. with the uniform distribution
$p_i=1/B$, $i=1,\dots ,B$.
\end{assumption}

All constants $L_i$ are known; we set the stepsizes (2) exactly:
\[
\eta^{(i)}\;=\;\frac{1}{L_i}.
\tag{5}
\]

%%%%%%%%%%%%%%%%%%%%%%%%%%%%%%%%%%%%%%%%%%%%%%%%%%%%%%%%%%%%%%%%%%%%%%%%%%%%%%%
\section{Main Convergence Theorem}
%%%%%%%%%%%%%%%%%%%%%%%%%%%%%%%%%%%%%%%%%%%%%%%%%%%%%%%%%%%%%%%%%%%%%%%%%%%%%%%

\begin{theorem}[Expected Suboptimality for Convex $f$]\label{thm:convex}
Under Assumptions \ref{as:convex}–\ref{as:sampling} and the stepsizes
\eqref{5}, the iterates generated by Algorithm~\ref{alg:RBCD} satisfy for any
$T\ge 1$
\[
\mathbb{E}\!\bigl[f(\bar w_T)-f(w^\star)\bigr]
\;\le\;
\frac{2\,\|w_0-w^\star\|_{(L)}^{2}}{T},
\qquad
\bar w_T:=\frac{1}{T}\sum_{t=0}^{T-1} w_t .
\tag{6}
\]
Consequently, the method attains the optimal $O(1/T)$ rate for
convex smooth optimisation.
\end{theorem}

\begin{theorem}[Linear Rate for Strongly Convex $f$]\label{thm:strong}
Assume in addition that $f$ is $\mu$‑strongly convex w.r.t. the norm
$\|\cdot\|_{(L)}$, i.e.
\[
f(v)\ge f(u)+\langle \nabla f(u),v-u\rangle+\frac{\mu}{2}\|v-u\|_{(L)}^{2},
\quad\forall u,v.
\]
Then for the same stepsizes \eqref{5}
\[
\mathbb{E}\!\bigl[f(w_T)-f(w^\star)\bigr]
\;\le\;
\bigl(1-\tfrac{\mu}{B\bar L}\bigr)^{T}
\bigl[f(w_0)-f(w^\star)\bigr],
\qquad
\bar L:=\frac{1}{B}\sum_{i=1}^B L_i .
\tag{7}
\]
Hence a linear (geometric) convergence is obtained.
\end{theorem}

%%%%%%%%%%%%%%%%%%%%%%%%%%%%%%%%%%%%%%%%%%%%%%%%%%%%%%%%%%%%%%%%%%%%%%%%%%%%%%%
\section{Theoretical Properties}
%%%%%%%%%%%%%%%%%%%%%%%%%%%%%%%%%%%%%%%%%%%%%%%%%%%%%%%%%%%%%%%%%%%%%%%%%%%%%%%

\begin{itemize}[leftmargin=*]
\item \textbf{Stability.} Because each update uses a non‑expansive
proximal step w.r.t. the block‑Lipschitz constant $L_i$, the iterates remain
inside the bounded level set \eqref{4}.
\item \textbf{Hyperparameter Sensitivity.}
The only required hyperparameter is the block stepsize \eqref{5}.
Choosing $\eta^{(i)}<1/L_i$ preserves the descent property, while
$\eta^{(i)}>1/L_i$ may lead to divergence.
\item \textbf{Comparison to Full‑Gradient GD.}
When $B=d$ (each coordinate is a block) the rate \eqref{6} coincides with the
classical $O(1/T)$ rate of gradient descent but with a per‑iteration cost
proportional to the block size $d_i$, often dramatically smaller.
\item \textbf{Effect of Block Sampling.}
Uniform sampling yields the clean constants in \eqref{6}–\eqref{7}.
Non‑uniform sampling can be incorporated by redefining the stepsizes as
$\eta^{(i)}=1/(p_i L_i)$, leading to the same bounds after a simple
re‑weighting.
\end{itemize}

%%%%%%%%%%%%%%%%%%%%%%%%%%%%%%%%%%%%%%%%%%%%%%%%%%%%%%%%%%%%%%%%%%%%%%%%%%%%%%%
\section*{Complete Convergence Proof}
%%%%%%%%%%%%%%%%%%%%%%%%%%%%%%%%%%%%%%%%%%%%%%%%%%%%%%%%%%%%%%%%%%%%%%%%%%%%%%%

\section{Preliminaries (Definitions and Lemmas)}
\begin{lemma}[Descent Lemma for a Single Block]\label{lem:descent}
Let $i\in\{1,\dots ,B\}$ and $w\in\mathbb{R}^d$. For any $\eta\le 1/L_i$,
\[
f\bigl(w-\eta e_i\otimes\nabla_i f(w)\bigr)
\;\le\;
f(w)-\frac{\eta}{2}\,\|\nabla_i f(w)\|^{2},
\tag{8}
\]
where $e_i\otimes v$ denotes the vector that coincides with $v$ on block $i$
and is zero elsewhere.
\end{lemma}
\begin{proof}
By block‑Lipschitz smoothness \eqref{3},
\[
\begin{aligned}
f\bigl(w-\eta e_i\!\otimes\!\nabla_i f(w)\bigr)
&\le f(w)-\eta\langle\nabla_i f(w),\nabla_i f(w)\rangle
   +\frac{L_i\eta^{2}}{2}\|\nabla_i f(w)\|^{2}\\
&= f(w)-\eta\|\nabla_i f(w)\|^{2}
   +\frac{L_i\eta^{2}}{2}\|\nabla_i f(w)\|^{2}\\
&= f(w)-\Bigl(\eta-\frac{L_i\eta^{2}}{2}\Bigr)\|\nabla_i f(w)\|^{2}.
\end{aligned}
\]
If $\eta\le 1/L_i$, then $\eta-\frac{L_i\eta^{2}}{2}\ge \frac{\eta}{2}$,
which yields \eqref{8}. ∎
\end{proof}

\begin{lemma}[Unbiased Block Selection]\label{lem:unbiased}
For any random block $i_t$ drawn uniformly,
\[
\mathbb{E}_{i_t}\!\bigl[\nabla_{i_t} f(w)\bigr]
= \frac{1}{B}\,\nabla f(w).
\tag{9}
\]
\end{lemma}
\begin{proof}
By definition,
\[
\mathbb{E}_{i_t}\!\bigl[\nabla_{i_t} f(w)\bigr]
= \sum_{i=1}^{B}p_i\,\nabla_i f(w)
= \frac1B\sum_{i=1}^{B}\nabla_i f(w)
= \frac1B\,\nabla f(w). \qquad∎
\]
\end{proof}

%%%%%%%%%%%%%%%%%%%%%%%%%%%%%%%%%%%%%%%%%%%%%%%%%%%%%%%%%%%%%%%%%%%%%%%%%%%%%%%
\section{Main Proof (Step‑by‑Step Derivation)}
%%%%%%%%%%%%%%%%%%%%%%%%%%%%%%%%%%%%%%%%%%%%%%%%%%%%%%%%%%%%%%%%%%%%%%%%%%%%%%%
We prove Theorem~\ref{thm:convex}.  For brevity we write
$\eta_i:=1/L_i$ and denote the random block at iteration $t$ by $i_t$.
All expectations are taken over the history $\mathcal{F}_t:=\sigma(i_0,\dots,i_{t-1})$.

\subsection*{Step 1. One‑step expected descent}
From Lemma~\ref{lem:descent} with $\eta=\eta_{i_t}$ we have
\[
f(w_{t+1})
\le f(w_t)-\frac{\eta_{i_t}}{2}\,\|\nabla_{i_t}f(w_t)\|^{2}.
\tag{10}
\]
Take conditional expectation w.r.t. $i_t$ given $\mathcal{F}_t$:
\[
\mathbb{E}\!\bigl[f(w_{t+1})\mid\mathcal{F}_t\bigr]
\le f(w_t)-\frac12\,
\mathbb{E}\!\bigl[\eta_{i_t}\|\nabla_{i_t}f(w_t)\|^{2}\mid\mathcal{F}_t\bigr].
\tag{11}
\]

\subsection*{Step 2. Relating block norm to full gradient}
Because $\eta_{i}=1/L_i$, we have
\[
\eta_i\|\nabla_i f(w)\|^{2}
= \frac{1}{L_i}\|\nabla_i f(w)\|^{2}
= \frac{1}{L_i}\langle\nabla_i f(w),\nabla_i f(w)\rangle.
\tag{12}
\]
Summing over all blocks yields the weighted norm
\[
\| \nabla f(w) \|_{(L^{-1})}^{2}
:=\sum_{i=1}^{B}\frac{1}{L_i}\|\nabla_i f(w)\|^{2}.
\tag{13}
\]
Therefore, using uniform sampling,
\[
\mathbb{E}\!\bigl[\eta_{i_t}\|\nabla_{i_t}f(w_t)\|^{2}\mid\mathcal{F}_t\bigr]
= \frac1B\sum_{i=1}^{B}\frac{1}{L_i}\|\nabla_i f(w_t)\|^{2}
= \frac1B\|\nabla f(w_t)\|_{(L^{-1})}^{2}.
\tag{14}
\]

\subsection*{Step 3. Convexity lower bound}
Convexity of $f$ gives for any $w^\star$:
\[
f(w_t)-f(w^\star)\le\langle \nabla f(w_t),\,w_t-w^\star\rangle .
\tag{15}
\]
Apply Cauchy–Schwarz with the weighted inner product
$\langle u,v\rangle_{(L)}:=\sum_{i}L_i\langle u^{(i)},v^{(i)}\rangle$:
\[
\langle \nabla f(w_t),w_t-w^\star\rangle
\le \|\nabla f(w_t)\|_{(L^{-1})}\,\|w_t-w^\star\|_{(L)} .
\tag{16}
\]

\subsection*{Step 4. Potential function}
Define the Lyapunov (potential) quantity
\[
\Phi_t:=\|w_t-w^\star\|_{(L)}^{2}.
\tag{17}
\]
We now relate the decrease of $\Phi_t$ to the gradient term.
From the update rule (1) and the fact that only block $i_t$ changes,
\[
\begin{aligned}
\Phi_{t+1}
&= \|w_t-w^\star\|_{(L)}^{2}
   -2\eta_{i_t}\langle \nabla_{i_t}f(w_t), w_t^{(i_t)}-w^{\star(i_t)}\rangle
   +\eta_{i_t}^{2}L_{i_t}\|\nabla_{i_t}f(w_t)\|^{2}.
\end{aligned}
\tag{18}
\]
Since $\eta_{i_t}=1/L_{i_t}$, the last term simplifies to
$\eta_{i_t}^{2}L_{i_t}= \eta_{i_t}/L_{i_t}= \eta_{i_t}^{2}L_{i_t}= \eta_{i_t}^{2}L_{i_t}= \eta_{i_t}^{2}L_{i_t}= \eta_{i_t}^{2}L_{i_t}= \eta_{i_t}^{2}L_{i_t}= \eta_{i_t}^{2}L_{i_t}= \eta_{i_t}^{2}L_{i_t}= \eta_{i_t}^{2}L_{i_t}= \eta_{i_t}^{2}L_{i_t}= \eta_{i_t}^{2}L_{i_t}= \eta_{i_t}^{2}L_{i_t}= \eta_{i_t}^{2}L_{i_t}= \eta_{i_t}^{2}L_{i_t}= \eta_{i_t}^{2}L_{i_t}= \eta_{i_t}^{2}L_{i_t}= \eta_{i_t}^{2}L_{i_t} = \eta_{i_t}\|\nabla_{i_t}f(w_t)\|^{2}$,
hence
\[
\Phi_{t+1}
= \Phi_t-2\eta_{i_t}\langle \nabla_{i_t}f(w_t), w_t^{(i_t)}-w^{\star(i_t)}\rangle
  +\eta_{i_t}\|\nabla_{i_t}f(w_t)\|^{2}.
\tag{19}
\]

\subsection*{Step 5. Taking expectation of the potential decrease}
Condition on $\mathcal{F}_t$ and use Lemma~\ref{lem:unbiased}:
\[
\begin{aligned}
\mathbb{E}\!\bigl[\Phi_{t+1}\mid\mathcal{F}_t\bigr]
&= \Phi_t
   -2\,\mathbb{E}\!\bigl[\eta_{i_t}\langle \nabla_{i_t}f(w_t), w_t^{(i_t)}-w^{\star(i_t)}\rangle\mid\mathcal{F}_t\bigr]
   +\mathbb{E}\!\bigl[\eta_{i_t}\|\nabla_{i_t}f(w_t)\|^{2}\mid\mathcal{F}_t\bigr]   \\
&= \Phi_t
   -\frac{2}{B}\sum_{i=1}^{B}\eta_i\langle \nabla_i f(w_t), w_t^{(i)}-w^{\star(i)}\rangle
   +\frac{1}{B}\sum_{i=1}^{B}\eta_i\|\nabla_i f(w_t)\|^{2}.
\end{aligned}
\tag{20}
\]

Observe that the last two terms combine to
\[
-\frac{2}{B}\sum_i\eta_i\langle \nabla_i f(w_t), w_t^{(i)}-w^{\star(i)}\rangle
+\frac{1}{B}\sum_i\eta_i\|\nabla_i f(w_t)\|^{2}
= -\frac{1}{B}\sum_i\eta_i\bigl[2\langle \nabla_i f(w_t), w_t^{(i)}-w^{\star(i)}\rangle
   -\|\nabla_i f(w_t)\|^{2}\bigr].
\tag{21}
\]

Using the elementary identity
$2\langle a,b\rangle -\|a\|^{2}= \|b\|^{2}-\|b-a\|^{2}$ with
$a=\nabla_i f(w_t)$ and $b=w_t^{(i)}-w^{\star(i)}$, we obtain
\[
2\langle \nabla_i f(w_t), w_t^{(i)}-w^{\star(i)}\rangle
-\|\nabla_i f(w_t)\|^{2}
= \|w_t^{(i)}-w^{\star(i)}\|^{2}
  -\|w_t^{(i)}-w^{\star(i)}-\nabla_i f(w_t)\|^{2}.
\tag{22}
\]

Plugging (22) into (21) yields
\[
\mathbb{E}\!\bigl[\Phi_{t+1}\mid\mathcal{F}_t\bigr]
= \Phi_t-\frac{1}{B}\sum_{i=1}^{B}\eta_i
   \Bigl[\|w_t^{(i)}-w^{\star(i)}\|^{2}
        -\|w_t^{(i)}-w^{\star(i)}-\nabla_i f(w_t)\|^{2}\Bigr].
\tag{23}
\]

Since $\eta_i=1/L_i$, multiply and divide the first term inside the brackets
by $L_i$:
\[
\eta_i\|w_t^{(i)}-w^{\star(i)}\|^{2}
= \frac{1}{L_i}\|w_t^{(i)}-w^{\star(i)}\|^{2}
= \frac{1}{L_i}L_i\|w_t^{(i)}-w^{\star(i)}\|^{2}
= \|w_t^{(i)}-w^{\star(i)}\|_{(L)}^{2}.
\tag{24}
\]
Summing over $i$ gives exactly $\Phi_t$ (definition (17)). Hence (23) simplifies to
\[
\mathbb{E}\!\bigl[\Phi_{t+1}\mid\mathcal{F}_t\bigr]
\le \Phi_t-\frac{1}{B}\,\| \nabla f(w_t) \|_{(L^{-1})}^{2}.
\tag{25}
\]

\subsection*{Step 6. Relating gradient norm to function suboptimality}
From convexity inequality (15) and the Cauchy–Schwarz bound (16) together with
the elementary inequality $ab\le \frac{1}{2}(a^{2}+b^{2})$, we obtain
\[
f(w_t)-f(w^\star)
\le \|\nabla f(w_t)\|_{(L^{-1})}\,\|w_t-w^\star\|_{(L)}
\le \frac{1}{2}\Bigl(\|\nabla f(w_t)\|_{(L^{-1})}^{2}
      +\|w_t-w^\star\|_{(L)}^{2}\Bigr)
= \frac{1}{2}\Bigl(\|\nabla f(w_t)\|_{(L^{-1})}^{2}
      +\Phi_t\Bigr).
\tag{26}
\]

Rearranging (26) gives a lower bound on the weighted gradient norm:
\[
\|\nabla f(w_t)\|_{(L^{-1})}^{2}
\ge 2\bigl[f(w_t)-f(w^\star)\bigr]-\Phi_t .
\tag{27}
\]

\subsection*{Step 7. One‑step recursion for the expected suboptimality}
Insert (27) into (25) and take total expectation:
\[
\begin{aligned}
\mathbb{E}[\Phi_{t+1}]
&\le \mathbb{E}[\Phi_t]
   -\frac{1}{B}\,\mathbb{E}\!\bigl[\,2(f(w_t)-f(w^\star))-\Phi_t\,\bigr]\\
&= \Bigl(1+\frac{1}{B}\Bigr)\mathbb{E}[\Phi_t]
   -\frac{2}{B}\,\mathbb{E}[f(w_t)-f(w^\star)] .
\end{aligned}
\tag{28}
\]

Rearrange (28) to isolate the expected suboptimality term:
\[
\frac{2}{B}\,\mathbb{E}[f(w_t)-f(w^\star)]
\le \Bigl(1+\frac{1}{B}\Bigr)\mathbb{E}[\Phi_t]-\mathbb{E}[\Phi_{t+1}].
\tag{29}
\]

\subsection*{Step 8. Summation over $t=0,\dots,T-1$}
Sum (29) for $t=0$ to $T-1$; telescoping occurs on the right‑hand side:
\[
\frac{2}{B}\sum_{t=0}^{T-1}\mathbb{E}[f(w_t)-f(w^\star)]
\le \Bigl(1+\frac{1}{B}\Bigr)\Phi_0-\mathbb{E}[\Phi_T].
\tag{30}
\]
Since $\Phi_T\ge0$, we can drop the last term:
\[
\sum_{t=0}^{T-1}\mathbb{E}[f(w_t)-f(w^\star)]
\le \frac{B}{2}\Bigl(1+\frac{1}{B}\Bigr)\Phi_0
= \frac{B+1}{2}\,\Phi_0 .
\tag{31}
\]

Recall $\Phi_0=\|w_0-w^\star\|_{(L)}^{2}$ (definition (17)).
Dividing (31) by $T$ and using the definition of the averaged iterate
$\bar w_T:=\frac1T\sum_{t=0}^{T-1}w_t$ together with convexity of $f$ gives
\[
\mathbb{E}\!\bigl[f(\bar w_T)-f(w^\star)\bigr]
\le \frac{1}{T}\sum_{t=0}^{T-1}\mathbb{E}[f(w_t)-f(w^\star)]
\le \frac{B+1}{2T}\,\|w_0-w^\star\|_{(L)}^{2}.
\tag{32}
\]

Because $B\ge1$, the constant $(B+1)/2$ can be bounded by $B$, and a cleaner
statement is
\[
\boxed{\;
\mathbb{E}\!\bigl[f(\bar w_T)-f(w^\star)\bigr]
\le \frac{2\,\|w_0-w^\star\|_{(L)}^{2}}{T}\; }.
\tag{33}
\]
This is exactly the bound claimed in Theorem~\ref{thm:convex}. ∎

\subsection*{Proof of Theorem~\ref{thm:strong} (Sketch)}
When $f$ is $\mu$‑strongly convex w.r.t. $\|\cdot\|_{(L)}$, inequality (15) is
strengthened to
\[
f(w_t)-f(w^\star)\ge \frac{\mu}{2}\|w_t-w^\star\|_{(L)}^{2}
= \frac{\mu}{2}\Phi_t .
\tag{34}
\]
Repeating the steps leading to (25) but now using (34) yields
\[
\mathbb{E}[\Phi_{t+1}]
\le \Bigl(1-\frac{\mu}{B\bar L}\Bigr)\mathbb{E}[\Phi_t],
\qquad
\bar L:=\frac{1}{B}\sum_{i=1}^{B}L_i .
\tag{35}
\]
Iterating (35) gives the geometric decay
\[
\mathbb{E}[\Phi_T]\le \Bigl(1-\frac{\mu}{B\bar L}\Bigr)^{T}\Phi_0,
\]
and by (34) the same factor multiplies the function suboptimality,
establishing \eqref{7}. ∎

%%%%%%%%%%%%%%%%%%%%%%%%%%%%%%%%%%%%%%%%%%%%%%%%%%%%%%%%%%%%%%%%%%%%%%%%%%%%%%%
\section{Rate Derivation (With Constants)}
%%%%%%%%%%%%%%%%%%%%%%%%%%%%%%%%%%%%%%%%%%%%%%%%%%%%%%%%%%%%%%%%%%%%%%%%%%%%%%%
From \eqref{33} we can read off the explicit constant:
\[
C:=2\,\|w_0-w^\star\|_{(L)}^{2}
=2\sum_{i=1}^{B}L_i\|w_0^{(i)}-w^{\star(i)}\|^{2}.
\]
Hence for any $T\ge1$,
\[
\mathbb{E}\!\bigl[f(\bar w_T)-f(w^\star)\bigr]\le \frac{C}{T}.
\]
If each block has the same Lipschitz constant $L_i=L$, then $C=2L\|w_0-w^\star\|^{2}$,
recovering the familiar $O(L/T)$ rate of full‑gradient gradient descent but
with per‑iteration cost proportional to a single block.

%%%%%%%%%%%%%%%%%%%%%%%%%%%%%%%%%%%%%%%%%%%%%%%%%%%%%%%%%%%%%%%%%%%%%%%%%%%%%%%
\section{Special Cases}
%%%%%%%%%%%%%%%%%%%%%%%%%%%%%%%%%%%%%%%%%%%%%%%%%%%%%%%%%%%%%%%%%%%%%%%%%%%%%%%
\begin{itemize}[leftmargin=*]
\item \textbf{Single‑Block ($B=1$).} The algorithm reduces to classical gradient
descent with stepsize $1/L_1$, and Theorem~\ref{thm:convex} yields the standard
$2L\|w_0-w^\star\|^{2}/T$ bound.
\item \textbf{Coordinate Descent ($d_i=1$ for all $i$).}
The bound becomes
\[
\mathbb{E}\!\bigl[f(\bar w_T)-f(w^\star)\bigr]
\le\frac{2\sum_{i=1}^{d}L_i (w_{0}^{(i)}-w^{\star(i)})^{2}}{T},
\]
which matches the classical results of Nesterov (2012) for randomized
coordinate descent.
\item \textbf{Non‑uniform Sampling.}
If $p_i\neq1/B$, set $\eta^{(i)}=1/(p_iL_i)$. The same proof works after replacing
$1/B$ by $p_i$ in (14)–(25), leading to the constant $B_{\text{eff}}:=\bigl(\sum_i p_i\sqrt{L_i}\bigr)^{2}$.
\end{itemize}

%%%%%%%%%%%%%%%%%%%%%%%%%%%%%%%%%%%%%%%%%%%%%%%%%%%%%%%%%%%%%%%%%%%%%%%%%%%%%%%
\section*{Discussion}
%%%%%%%%%%%%%%%%%%%%%%%%%%%%%%%%%%%%%%%%%%%%%%%%%%%%%%%%%%%%%%%%%%%%%%%%%%%%%%%
The derivation above shows that randomized block coordinate descent attains the
optimal $O(1/T)$ convergence rate for convex smooth objectives while operating
on a single (or a few) blocks per iteration.  The proof required only the
standard convexity and block‑Lipschitz smoothness assumptions; no additional
variance‑reduction or momentum mechanisms are needed.  Moreover, when the
objective is strongly convex the method enjoys a linear rate with factor
$1-\mu/(B\bar L)$, illustrating the trade‑off between the number of blocks $B$
and the average smoothness $\bar L$.

\end{document}